\documentclass[10pt,a4paper]{amsart}
\usepackage[margin=19mm]{geometry}
\usepackage{amsmath,amssymb,mathtools}
\usepackage[scr=rsfs,cal=euler]{mathalfa}
\usepackage{xfrac}
\usepackage[unicode,hidelinks,bookmarks=false]{hyperref}
\newcommand{\ie}{i.e.\ }
\newcommand{\cf}{cf.\ }
\newcommand{\resp}{respectively\ }
\newcommand{\Subsection}{\subsection}
\providecommand{\nobreakdash}{\nobreak\hbox{-}}
\setlength{\emergencystretch}{2em}
\multlinegap0pt
\usepackage{mathtools}
\usepackage{amssymb}
\usepackage[shortlabels]{enumitem}
\usepackage{xcolor}
\newtheorem{theorem}{Theorem}
\title{From curvature to Kovacic: a geometric approach to integrability of scalar ODEs}
\author{A. J. Pan-Collantes, J. A. \'Alvarez-Garc\'ia}
\date{Source: arXiv:2604.05752v1 (CC BY 4.0)}
\begin{document}
\maketitle

\noindent\emph{Source and licence.} From curvature to Kovacic: a geometric approach to integrability of scalar ODEs. Version arXiv:2604.05752v1 (CC BY 4.0), distributed by arXiv under the Creative Commons Attribution 4.0 International licence, http://creativecommons.org/licenses/by/4.0/, as recorded in the arXiv licence metadata for this exact version and shown on the abstract page https://arxiv.org/abs/2604.05752v1. Full licence notice: \texttt{licenses/arxiv-2604.05752-CC-BY-4.0.txt}.\par\smallskip

\noindent\textit{Editorial scope.} Counted unit: the article's theorem thm:riccati\_connection (source0.tex lines 184--186). The article's complete original proof of it is delivered with it (source0.tex lines 188--211, one \texttt{proof} environment of 24 lines). No mathematics is written, repaired or re-derived: the statement and the proof are the article's own lines, in the article's own notation, and no retained line is substituted or re-flowed.  Retained uncounted as the source-local context the proof needs: lines 104--106; lines 120--122; lines 136--138; lines 153--155.  Nothing was omitted from any retained span, so the package carries no omission record.  No retained line contains a citation, so the wrapper carries no bibliography and no bibliography entry.  No retained line contains a figure environment, an image-inclusion command or drawing code, so no image is delivered and the package declares no asset dependency.  Editorial formatting only, in addition: an AMS \texttt{amsart} preamble that restates the article's own declarations and macros used by the retained lines, namely the environments \texttt{theorem} on separate counters, together with the standard packages those lines need.  The retained environments print in this document as Theorem 1 in order of appearance, whereas the article prints the same items with the numbering of its own full document.  The complete native source of the article as distributed in the arXiv e-print for this exact version is bundled unchanged as \texttt{sources/197977/source0.tex}.
\begin{equation} \label{eq:ode1}
    u'(x) = \phi(x, u)
\end{equation}

\begin{equation} \label{eq:vector_field}
    A = \partial_x + \phi\partial_u,
\end{equation}

\begin{equation} \label{eq:curvature_formula}
    \mathcal{K}(x,u) = -\partial_u(A(\phi)),
\end{equation}

\begin{equation}\label{Riccatieq}
p'(x) + p(x)^2 + \kappa(x) = 0
\end{equation}

\begin{theorem}\label{thm:riccati_connection}
The curvature corresponding to equation \eqref{eq:ode1} satisfies $\mathcal{K}(x,u)=\kappa(x)$ on $D$ if and only if, for every solution $f$, the function $p_f(x):=\phi_u(x,f(x))$ satisfies the Riccati equation \eqref{Riccatieq} on its interval of definition.
\end{theorem}

\begin{proof}%nueva
Let $f$ be any solution of \eqref{eq:ode1}. By the chain rule and $f'(x)=\phi(x,f(x))$,
\[
p_f'(x)=\frac{d}{dx}\phi_u(x,f(x))=\phi_{xu}(x,f(x))+\phi(x,f(x))\,\phi_{uu}(x,f(x))=A(\phi_u)(x,f(x)),
\]
where $A$ is the associated vector field \eqref{eq:vector_field}.


By equation \eqref{eq:curvature_formula} we have
\[
\mathcal{K}(x,u)=-\partial_u\bigl(A(\phi)\bigr)(x,u)=-(\phi_{xu}+\phi_u^{2}+\phi\,\phi_{uu})(x,u)=-A(\phi_u)(x,u)-\phi_u^2(x,u). %% simplified notation
\]
Evaluating at $u=f(x)$ and using the previous identities yields, along the solution curve,
\begin{equation}\label{eq:riccati_identity}
p_f'(x)+p_f(x)^2+\mathcal{K}(x,f(x))=0.
\end{equation}
If $\mathcal{K}(x,u)=\kappa(x)$, this reduces to the Riccati equation \eqref{Riccatieq} for $p_f$, for every solution $f$.

Conversely, assume that there exists a function $\kappa=\kappa(x)$ such that for every solution $f$ the corresponding $p_f$ satisfies \eqref{Riccatieq}. Comparing with the identity \eqref{eq:riccati_identity} gives $\mathcal{K}(x,f(x))=\kappa(x)$ along each solution curve. Given any point $(x_0,u_0)$ in the domain of $\phi$, local existence for the initial value problem $u'=\phi(x,u)$, $u(x_0)=u_0$, provides a solution $f$ with $f(x_0)=u_0$, hence
\[
\mathcal{K}(x_0,u_0)=\mathcal{K}(x_0,f(x_0))=\kappa(x_0).
\]
Since $(x_0,u_0)$ is arbitrary, it follows that $\mathcal{K}(x,u)=\kappa(x)$ on the domain.
\end{proof}

\end{document}
